It remains to prove the last assertion, and one may clearly suppose
$G=G^0$. For Cartan subgroups, it follows from the analogous assertion
in $G'$ (BIBLE 6 th. 4 cor. 4) and from XII 6.6 e). For maximal tori,
it follows from the preceding assertion, since by XII 6.6 c) the
maximal tori of a smooth algebraic group are the maximal tori of
its Cartan subgroups.

\begin{corollary}{4.3}\label{XIV.4.3}
Suppose $G$ connected. Then every element of $G$ is contained in a
Borel subgroup of $G$.
\end{corollary}
One is reduced to the same statement in $G'$, which is well known
(BIBLE 6 th. 5 d)).

\begin{corollary}{4.4}\label{XIV.4.4}
Let $B$ be a Borel subgroup of $G$, $C$ a Cartan subgroup of $B$,
and $N$ its normalizer in $G$; then $N\cap B=C$.
\end{corollary}
Indeed, $N\cap B=\Norm_B(C)$, so one is reduced to showing that
when $G$ is connected and solvable, a Cartan subgroup $C$ is its
own connected normalizer. Now, with the preceding notation,
$C$ is the inverse image of a Cartan subgroup $C'$ of $G'$, so
one is reduced to the case where $G$ is affine. Since one knows
that the normalizer of a Cartan subgroup is smooth ($C$ being
its own connected normalizer; cf., for example, XII 6.6 c)),
it suffices to see that $C$ and $N$ have the same points with
values in $k$, which is precisely BIBLE th. 6 d).
% typo?: kept as printed: connected normalizer in the reduction.

\begin{definition}{4.5}\label{XIV.4.5}
Let $G$ be a smooth group prescheme of finite presentation over
a prescheme $S$. One calls a Borel subgroup of $G$ any smooth
group subprescheme $B$ of finite presentation of $G$ such that
for every $s\in S$, the geometric fiber $B_{\bar s}$ is a Borel
subgroup of $G_{\bar s}$.
\end{definition}
Thus this is, as one verifies at once, a notion stable under
base change and local for the fpqc topology (for if $k'$ is
an algebraically closed extension of an algebraically closed
field $k$, then an algebraic subgroup $B$ of the smooth
algebraic group $G$ over $k$ is a Borel subgroup of $G$ if
and only if $B_{k'}$ is one of $G_{k'}$). It follows from this
definition that if $G$ is a smooth algebraic group over an
arbitrary field $k$, and $B$ a Borel subgroup of $G$, then
$G/B$ is a projective variety, every maximal torus $T$ of
$B$ is a maximal torus of $G$, its normalizer in $B$ is
identical to its centralizer $C$, and is a Cartan subgroup
of $G$ when $G$ is connected.

\begin{remarks}{4.6}\label{XIV.4.6}
Unfortunately, it is no longer true in general (even if $G$
is affine over $S$ and $S$ is the spectrum of the algebra
of dual numbers over an algebraically closed field $k$)
that two Borel subgroups of $G$ are conjugate locally for
the fpqc topology. As a consequence of this regrettable
fact, let us point out that if $G$ is a smooth affine
connected algebraic group over an imperfect field $k$,
it is not in general possible to define naturally a
homogeneous space $D$ under $G$, playing the role of a
flag variety, i.e. of the variety of Borel subgroups of
$G$ (which over the algebraic closure $\bar k$ of $k$
would therefore be isomorphic to $G_{\bar k}/\bar B$,
where $\bar B$ is a Borel subgroup of $G_{\bar k}$).
Indeed, when the quotient of $G_{\bar k}$ by its radical
$\bar R$ is an adjoint semisimple group, then the kernel
of $G_{\bar k}\to\underline{\Aut}_{\bar k}(\bar D)$
is the radical of $G_{\bar k}$; hence if $\bar D$ comes
from a homogeneous space $D$ under $G$, the radical
$\bar R$ comes from a subgroup $R$ of $G$. Now one
easily constructs examples where $G_{\bar k}/\bar R$
is adjoint but $\bar R$ is not ``defined over $k$''.
It is easy to see that under these conditions, the
functor $\mathscr B:(\Sch)^\circ_{/S}\to\Ens$ such
that $\mathscr B(S')=$ the set of Borel subgroups of
$G_{S'}$ is not representable by a smooth
$S$-prescheme. From the infinitesimal point of view
(III \S 3), the failure of the conjugacy theorem is
expressed by the fact that if $B$ is a Borel subgroup
of the smooth algebraic group $G$, the cohomology group
$\mathrm H^1(B,\mathfrak g/\mathfrak b)$\nde{Where
$\mathfrak b$ is the Lie algebra of $B$.} may be nonzero.

On the other hand, we shall see in a later expos\'e
that when $G$ is semisimple, or more generally reductive,
such unpleasant phenomena do not occur. It is doubtless
these phenomena, as well as the absence of good existence
theorems, which explain why Borel subgroups play only
a relatively unobtrusive role in the study of general
group schemes from the scheme-theoretic point of view,
whereas they will dominate the theory of semisimple
group schemes in later expos\'es.
\end{remarks}

\begin{proposition}{4.8}\label{XIV.4.8}
\nde{There is no number 4.7.} Let $G$ be a smooth
$S$-group prescheme of finite presentation with
connected fibers, and $B$ a Borel subgroup of
$G$; then $B$ is identical to its own normalizer,
and it is a closed subprescheme of $G$.
\end{proposition}
Indeed, by XII 7.10, one is reduced to proving
that over an algebraically closed field $k$,
every element of $G(k)$ which normalizes $B$
is in $B(k)$, which for affine $G$ is a
fundamental result of Chevalley (BIBLE 9 th. 1);
the general case reduces to it at once by
the reduction already used in 4.2.

\begin{remarks}{4.8.1}\label{XIV.4.8.1}
One may generalize definition 4.5 by introducing
also the notion of parabolic subgroup of $G$:
one thus calls a group subprescheme $P$ of
$G$, smooth and of finite presentation over
$S$, such that for every $s\in S$, the
geometric fiber $P_{\bar s}$ is a parabolic
subgroup of $G_{\bar s}$, i.e. contains a
Borel subgroup of $G_{\bar s}$. Proposition
4.8 extends (with the same proof of reduction
to the ``set-theoretic'' statement, which is
known) to the case of a parabolic subgroup
$P$ of $G$. Note the following consequence
of this result (cf. XVI). If $P$ is a
parabolic subgroup of $G$, then $G/P$ is
representable by a quasi-projective prescheme
of finite presentation over $S$ (N.B. one
assumes $G$ has connected fibers). Moreover,
$G/P$ is clearly smooth over $S$, and
moreover has connected proper geometric
fibers, whence one may easily conclude,
using EGA III 5.5.1, that $D=G/P$ is in
fact proper, hence projective, over $S$.
Moreover, if its relative dimension is
$n$, it is known that the invertible
sheaf $\Omega^n_{D/S}$ is such that its
inverse induces ample sheaves on the
geometric fibers of $D/S$, hence
(EGA III 4.7.1) $(\Omega^n_{D/S})^{-1}$
is ample on $D$ relative to $S$.
\end{remarks}

One easily sees, by reduction to the
affine case and to the case of an
algebraically closed base field, that
if $u:G\to G'$ is an epimorphism of
smooth algebraic groups, then for
every Borel subgroup $B$ of $G$,
$u(B)=B'$ is a Borel subgroup of $G'$.
We are interested in the case where
one thus obtains a bijective
correspondence between Borel subgroups
of $G$ and of $G'$:

\begin{proposition}{4.9}\label{XIV.4.9}
Let $G,G'$ be two smooth $S$-group
preschemes of finite presentation
with connected fibers, and
$u:G\to G'$ a faithfully flat
homomorphism of groups, i.e.
surjective\nde{Add a reference here?}.
One assumes that one is in one
of the following two cases
(where one has set $N=\Ker u$):

\noindent a) $N$ is central in $G$.

\noindent b) $S$ is the spectrum
of a field $k$, and if $\bar k$
denotes an algebraic closure of
it, $N_{\bar k}=\Ker u_{\bar k}$
is contained in the radical of
$G_{\bar k}$, i.e. in the
largest smooth connected solvable
invariant subgroup of $G_{\bar k}$.

Then the map $B'\mto u^{-1}(B')$
induces a bijection from the set
of Borel subgroups of $G'$ onto
the analogous set for $G$.
\end{proposition}
Case b) follows at once from the
correspondence between algebraic
subgroups of $G'$ and algebraic
subgroups of $G$ containing $N$,
and the fact that when $k$ is
algebraically closed, the Borel
subgroups of $G$ contain the
radical of $G$ (which is immediate
by the argument preceding 4.2).

Let us prove case a). By XII
7.12, the map $H'\mto H=u^{-1}(H')$
establishes a one-to-one
correspondence between group
subpreschemes $H'$ of $G'$
which are smooth and of finite
presentation over $S$, with
connected fibers, and which
have at every $s\in S$ the
same reductive rank and the
same nilpotent rank as $G'$,
and group subpreschemes $H$
of $G$ having the analogous
properties. Now the Borel
subgroups (of $G'$, or of
$G$) have the properties in
question. It remains to
prove that if $H',H$
correspond, then $H'$ is
a Borel subgroup of $G'$
if and only if $H$ is
one in $G$. By definition,
this question reduces to
the case where $S$ is the
spectrum of an algebraically
closed field. Now, since
$N$ is central in $G$,
hence in $H$, it follows
that $H$ is solvable if
and only if $H'$ is.
Finally, taking into account
the correspondence between
algebraic subgroups of $G'$
and algebraic subgroups of
$G$ containing $N$, one sees
at once that $H'$ possesses
the maximality property of
definition 4.1 if and only
if $H$ possesses it, which
completes the proof.

\begin{corollary}{4.10}\label{XIV.4.10}
With the notation of 4.7,
if $B'$ and $B$ are Borel
subgroups of $G'$ and $G$
which correspond, one has
\[
 \mathfrak b=\Lie(u)^{-1}(\mathfrak b')
\]
where $\mathfrak g,\mathfrak g',\mathfrak b,\mathfrak b'$
are the Lie algebras of
$G,G',B,B'$, and where
$\Lie(u):\mathfrak g\to\mathfrak g'$
is the homomorphism deduced
from $u$.
% typo?: kept as printed: reference 4.7, which does not exist.
\end{corollary}
This statement follows
trivially from the definitions
and the relation $u^{-1}(B')=B$.

We can now prove the main
result of the present section:
\begin{theorem}{4.11}\label{XIV.4.11}
Let $G$ be a smooth
algebraic group over an
algebraically closed field
$k$, and $\mathfrak g$
its Lie algebra. Then
$\mathfrak g$ is equal
to the union of the
Lie algebras $\mathfrak b$
of the Borel subgroups
$B$ of $G$.
\end{theorem}
One may clearly suppose
$G$ connected. Let $R$
be the radical of $G$,
and let $G'=G/R$.
Then 4.9 b) and 4.10
reduce one to proving
theorem 4.11 for $G'$
in place of $G$, i.e.
one may suppose $G$
semisimple. Let then
$Z$ be the center of
$G$, identical to the
reductive center, and
let $G'=G/Z$.
The same argument
(now using 4.9 a))
reduces one to proving
the theorem for $G'$,
i.e. one may suppose
$G$ adjoint semisimple.
Let $B$ be a Borel
of $G$, $T$ a maximal
torus of $B$, hence
of $G$, and
$\mathfrak b,\mathfrak t$
the Lie algebras.
By 3.18, $T$ is a
subgroup of type (C)
of $G$, i.e.
$\mathfrak t$ is a
Cartan subalgebra of
$\mathfrak g$, so
the union of the
conjugates of
$\mathfrak t$ is
dense in $\mathfrak g$
(XIII 5.1 (i)
$\Rightarrow$ (vii)).
A fortiori the union
of the conjugates of
$\mathfrak b$ is
dense in $\mathfrak g$.
Now let $X$ be the
closed subscheme of
$G/B\times\mathrm W(\mathfrak g)$
whose points with
values in $k$ are
the $(g',x)$ such
that $x\in\Ad(g)\cdot\mathfrak b$
(cf. XIII 1).
Then the morphism
$\psi:X\to\mathrm W(\mathfrak g)$
induced by the
second projection
is proper since
$G/B$ is proper
over $k$; on the
other hand one has
just seen that it
is dominant, so
it is surjective,
which proves 4.11.

The only result
of the present
section that we
shall use in the
rest of this
expos\'e is the
following corollary:
\begin{corollary}{4.12}\label{XIV.4.12}
Let $k$ be an
infinite field,
$G$ a smooth
algebraic group
over $k$, $T$
a maximal torus
of $G$,
$\mathfrak g$ and
$\mathfrak t$ the
Lie algebras,
and $u:G\to\GL(V)$
a linear representation
($V$ a finite-dimensional
vector space over
$k$), whence a
representation of
Lie algebras
$u':\mathfrak g\to\mathfrak{gl}(V)$
making $V$ a
$\mathfrak g$-module.
Then the minimum
of the nullity
of $u'(x)$
($x\in\mathfrak g$)
is attained for
an element
$x\in\mathfrak t$.
\end{corollary}
One is reduced
at once to the
case where $k$
is algebraically
closed. One may
clearly suppose
$G$ connected,
and, dividing
$G$ by
$(\Ker u)_{\mathrm{red}}$,
one may suppose
$G$ affine.
Using 4.11 and
the conjugacy
theorem for
maximal tori of
$G$, one is
reduced to the
case where $G$
is moreover
solvable. Then
$G$ is a
semidirect product
$T\cdot V$,
where $V$ is
the ``unipotent
part'' of $G$,
which is a
smooth connected
unipotent group
(BIBLE 6 th. 3).
Thus $\mathfrak g$
is a direct sum
(as a vector
space) of the
Lie subalgebras
$\mathfrak t$ and
$\mathfrak v=\Lie(V)$.
By the Lie--Kolchin
theorem (BIBLE 6
th. 1),
$\mathfrak v$ admits
a composition
series by stable
subspaces
$\mathfrak v_i$
such that
$\mathfrak v_i/\mathfrak v_{i+1}=\mathfrak w_i$
has dimension $1$.
Then for each
$i$, one has
an induced
representation
$u_i:G\to\GL(\mathfrak w_i)=\Gm$
and the
corresponding
homomorphism of
Lie algebras
$u'_i:\mathfrak g\to k$,
so that for
every
$x\in\mathfrak g$,
the nullity of
$u'(x)$ is
equal to the
number of $i$
such that
$u'_i(x)=0$.
Since $V$ is
unipotent, the
$u_i$ are
trivial on $V$,
hence the
$u'_i$ are
trivial on
$\mathfrak v$,
which proves
that if $x=t+v$
($t\in\mathfrak t$,
$v\in\mathfrak v$),
then $u'_i(x)=u'_i(t)$
for every $i$,
hence the
nullity of
$u'(x)$ is
equal to that
of $u'(t)$.
Assertion 4.12
follows at once.
% typo?: kept as printed: V names the unipotent group, and v is the representation space here.

\sgasection{5}{Relations between Cartan subgroups and Cartan subalgebras}
Applying 4.12
to the adjoint
representation of
$G$, one finds:
\begin{theorem}{5.1}\label{XIV.5.1}
Let $G$ be a
smooth algebraic
group over an
infinite field,
$T$ a maximal
torus of $G$,
and
$\mathfrak g\supset\mathfrak t$
the Lie
algebras; then
$\mathfrak t$
contains a
regular element
of $\mathfrak g$.
\end{theorem}

\begin{corollary}{5.2}\label{XIV.5.2}
Let $G$ be a
smooth $S$-group
prescheme of
finite presentation,
and $\mathfrak g$
its Lie algebra.

\noindent a) Conditions
$(\mathrm C_0)$ to
$(\mathrm C_2)$ of
2.9 on
$\mathfrak g$ are
equivalent; in
particular, if
the infinitesimal
rank of the
fibers of
$\mathfrak g$ at
the points of
$S$ is locally
constant, then
$\mathfrak g$
admits a Cartan
subalgebra locally
for the \'etale
topology, hence
(by 3.9 a))
$G$ admits a
subgroup of type
(C) locally for
the \'etale
topology.

\noindent b) Let
$H$ be a group
subprescheme of
$G$ smooth and
of finite
presentation over
$S$, with
connected fibers
having the same
reductive rank
as $G$ at each
$s\in S$ (for
example, $H$ is
a maximal torus
or a Cartan
subgroup of $G$),
and $D$ a
subgroup of type
(C) of $G$;
then one has
$H\subset D$
if and only if
one has
$\mathfrak h\subset\mathfrak d$.

\noindent c) Suppose
condition
$(\mathrm C_0)$
holds, i.e. the
infinitesimal rank
of $G$ is locally
constant. Let $H$
be a group
subprescheme of
$G$ smooth and
of finite
presentation over
$S$, with
nilpotent connected
fibers, having
at every point
the same reductive
rank as $G$
(for example,
$H$ is a maximal
torus or a Cartan
subgroup of $G$).
Then $H$ is
contained in one
and only one
subgroup $D$ of
type (C) of $G$.

\noindent d) Suppose
that $G$ admits
a Cartan subgroup
locally for fpqc;
then so does
every subgroup
$D$ of type
(C) of $G$.
\end{corollary}
\begin{proof}
a) Suppose condition
$(\mathrm C_0)$
holds, and let
us prove that
the same holds
for $(\mathrm C_2)$,
i.e. that every
quasi-regular
section $a$ of
$\mathfrak g$ is
regular. One
reduces as usual
to the case
of $S$ affine
noetherian, then,
the question being
``infinitesimal''
(2.15 b)), to
the case where
$S$ is artinian
local (then
$(\mathrm C_0)$
is moreover
trivially satisfied).
One may suppose
in addition that
the residue field
$k$ of $S$ is
infinite. Note
that by 3.7
it suffices to
establish that
$\mathfrak g$
admits a Cartan
subalgebra. Let
$T$ be a maximal
torus of $G$;
by 5.1 there
exists a
quasi-regular
element $t$
contained in the
Lie algebra
$\mathfrak t$ of
$T$; let us
prove that it
is regular,
which will finish
the proof.
Consider the
linear representation
of $T$ in
$\mathfrak g$
induced by the
adjoint representation
of $G$; there
thus exists a
finite set
$(u_i)_{i\in I}$
of characters of
$T$ such that
$\mathfrak g$
decomposes as a
direct sum of
submodules
$\mathfrak g_i$
stable under $T$,
$T$ acting on
$\mathfrak g_i$
by $u_i$ (cf.
I \S 4.7.3).
Let
$u'_i:\mathfrak t\to A$
be the homomorphism
of Lie algebras
deduced from
$u_i:T\to\Gm$
(N.B. $A$
denotes the ring
of $S$).
Consider the
homomorphisms
$u_{i0}$ and
$u'_{i0}$ deduced
from the preceding
ones by passage
to the fibers,
i.e. by the
base change
$A\to k$.
Let $I'$ be
the set of
$i\in I$ such
that $u'_{i0}\ne0$,
and let
$I''=I-I'$.
The fact that
$t$ is regular
is expressed by
the condition
$u'_{i0}(t_0)\ne0$
for every
$i\in I'$,
hence $u'_i(t)$
invertible for
$i\in I'$.
The Cartan
subalgebra of
$\mathfrak g_0$
defined by $t$,
i.e. the kernel
subspace of the
semisimple endomorphism
$\ad(t_0)$ of
$\mathfrak g_0$,
is
$\sum_{i\in I''}(\mathfrak g_i)_0$.
Consider
\[
 \mathfrak d=\sum_{i\in I''}\mathfrak g_i,
\]
then $\ad(t)$
is nilpotent on
$\mathfrak d$;
on the other
hand, it is
an automorphism
of
$\mathfrak g/\mathfrak d\simeq\sum_{i\in I'}\mathfrak g_i$.
By 2.6, $t$
is therefore regular.

\noindent b) Follows
from 3.12 a),
$H$ satisfying
the hypothesis
that every
geometric fiber
of $\mathfrak h$
contains a regular
element of that
of $\mathfrak g$,
by 5.1.

\noindent c) By b),
one is reduced
to proving that
$\mathfrak h$ is
contained in one
and only one
Cartan subalgebra
of $\mathfrak d$.
By a), one
knows moreover
that $\mathfrak g$
satisfies
$(\mathrm C_2)$.
One is therefore
reduced to the
% typo?: kept as printed: a Cartan subalgebra of d.

\begin{lemma}{5.3}\label{XIV.5.3}
Let $\mathfrak g$
be a Lie algebra
over $S$ which
is a locally
free module of
finite type and
satisfies condition
$(\mathrm C_2)$
of 2.9. Let
$\mathfrak h$ be
a Lie subalgebra
of $\mathfrak g$
satisfying the
following conditions:
it is a module
locally a direct
summand, it is
locally nilpotent,
and for every
$s\in S$ the
geometric fiber
$\mathfrak h_{\bar s}$
contains a regular
element of
$\mathfrak g_{\bar s}$.
Then $\mathfrak h$
is contained in
one and only one
Cartan subalgebra
of $\mathfrak g$.
\end{lemma}
(N.B. in the
case of interest
to us,
$\mathfrak h$
satisfies the
stated conditions:
it is locally
nilpotent for
$H$ has nilpotent
fibers, and the
condition on regular
elements follows
from 5.1.) Since
5.3 is local
for the fpqc
topology, it
suffices to
prove that at
a point
$s\in S$ such
that $\kappa(s)$
is infinite,
there exists an
open neighborhood
$U$ of $s$
such that
existence and
uniqueness hold
for every base
change $S'\to S$
factoring through
$U$. Take a
regular element
of
$\mathfrak g\otimes\kappa(s)$
contained in
$\mathfrak h\otimes\kappa(s)$;
extend it to
a section $a$
of $\mathfrak h$
over an open
neighborhood $U$
of $s$; by
$(\mathrm C_2)$
one may suppose
that this section
is regular (2.10).
One may suppose
$U=S$. A Cartan
subalgebra of
$\mathfrak g$
containing
$\mathfrak h$
contains $a$,
hence is identical
to
$\mathfrak d=\operatorname{Nil}(a,\mathfrak g)$
(2.6), whence
uniqueness.
Moreover, since
$\mathfrak h$ is
locally nilpotent,
one indeed has
$\mathfrak h\subset\mathfrak d$,
which proves
existence.

\noindent d) This
is a particular
case of XII
7.9 d).
\end{proof}

\begin{corollary}{5.4}\label{XIV.5.4}
Let $G$ be a
smooth connected
algebraic group
over an
algebraically
closed field $k$,
$H$ a connected
algebraic subgroup
such that
$\mathfrak h$
contains a Cartan
subalgebra of
$\mathfrak g$,
i.e. such that
$H$ contains a
subgroup of type
(C) of $G$.
Then the number
of conjugates
of $H$ containing
a regular element
of $G(k)$ is
equal to the
number of conjugates
of $\mathfrak h$
containing a regular
element of
$\mathfrak g$.
% typo?: kept as printed: H is not assumed smooth.
\end{corollary}
Indeed, let $g$
be a regular
element of $G(k)$,
and $C$ the
unique Cartan
subgroup of $G$
containing $g$
(XIII 2); then
the conjugates
of $H$ containing
$g$ are those
containing $C$
(XIII 2.8 b)).
Likewise, let
$x$ be a regular
element of
$\mathfrak g$;
then if a
conjugate
$\mathfrak h'$
of $\mathfrak h$
contains $x$,
then
$\ad(x)_{\mathfrak g/\mathfrak h'}$
is injective
(XIII 5.4), so
$\mathfrak h'$
contains
$\operatorname{Nil}(x,\mathfrak g)=\mathfrak d$;
hence the number
of conjugates
of $\mathfrak h$
containing $x$
is equal to the
number of conjugates
containing the
Cartan subalgebra
$\mathfrak d$.
Moreover,
$\mathfrak d$ is
the Lie algebra
of a subgroup
$D$ of type (C)
of $G$.
One may clearly
suppose $C\subset D$,
and assertion
5.4 will follow
from this:
for $H$ to
contain $C$,
it is necessary
and sufficient
that $\mathfrak h$
contain
$\mathfrak d$.
Indeed, by XIII
5.5 the relation
$\mathfrak h\supset\mathfrak d$
implies
$H\supset D$,
and a fortiori
$H\supset C$.
Conversely, by
5.1,
$\mathfrak t$
contains a regular
element $x$ of
$\mathfrak g$;
then $H\supset C$
implies
$\mathfrak h\ni x$,
hence, as already
noted, this implies
$\mathfrak h\supset\mathfrak d$.
This completes
the proof.

Let $G$ be
as in 5.2,
and suppose
that the
infinitesimal
rank of the
fibers of $G$
remains locally
constant
(condition
$(\mathrm C_0)$).
Then by 5.2 c),
one finds a
homomorphism of
functors on
$(\Sch)^\circ_{/S}$:
\[
 \mathscr C\longrightarrow\mathscr D
\]
where
\[
\begin{aligned}
 \mathscr C(S')&=\text{set of Cartan subgroups of }G_{S'},\\
 \mathscr D(S')&=\text{set of subgroups of type (C) of }G_{S'}.
\end{aligned}
\]
